\documentclass[11pt,a4paper]{article}
\usepackage[a4paper,margin=2cm,headheight=14pt]{geometry}
\usepackage{fontspec}
\setmainfont{Arial}
\setmonofont{Courier New}
\usepackage{polyglossia}
\setdefaultlanguage{brazil}
\usepackage{xcolor,tikz,listings,hyperref,fancyhdr,enumitem}
\usetikzlibrary{arrows.meta,positioning}
\definecolor{azul}{HTML}{174A72}\definecolor{azulclaro}{HTML}{EAF2F8}\definecolor{cinza}{HTML}{F1F4F7}\definecolor{verde}{HTML}{26734D}\definecolor{ouro}{HTML}{D59B32}
\hypersetup{colorlinks=true,linkcolor=azul,urlcolor=azul}
\pagestyle{fancy}\fancyhf{}\lhead{Integração Flask + SUAP}\rhead{Guia de implementação}\cfoot{\thepage}
\lstset{basicstyle=\ttfamily\small,breaklines=true,frame=single,backgroundcolor=\color{cinza},columns=fullflexible,showstringspaces=false,keywordstyle=\color{azul}\bfseries,commentstyle=\color{verde}}
\begin{document}
\begin{titlepage}
\thispagestyle{empty}
\pagecolor{azul}
\vspace*{1.25cm}
{\color{white}\Large DESENVOLVIMENTO WEB}\par\vspace{1.35cm}
{\color{white}\Huge\bfseries Tutorial de integração\\[.25cm]com o SUAP}\par\vspace{.45cm}
{\color{white}\Large Login OAuth 2.0 e consumo da API\\em uma aplicação Flask}\par
\vspace{.8cm}{\color{white}\large Guia prático do cadastro ao primeiro login}\par
\vspace{4.25cm}
\begin{center}
\begin{tikzpicture}[node distance=1.2cm,>=Stealth]
\node[draw=azul,fill=azulclaro,rounded corners=4pt,minimum width=3.2cm,minimum height=1.05cm,align=center] (u) {Pessoa usuária};
\node[draw=azul,fill=azulclaro,rounded corners=4pt,minimum width=3.2cm,minimum height=1.05cm,align=center,right=of u] (f) {Aplicação Flask};
\node[draw=verde,fill=green!8,rounded corners=4pt,minimum width=3.2cm,minimum height=1.05cm,align=center,right=of f] (s) {SUAP / OAuth 2.0};
\draw[-{Stealth},thick,white] (u)--node[above=4pt,draw=none,font=\small,text=white]{1. iniciar}(f);
\draw[-{Stealth},thick,ouro,bend left=24] (f.north east) to node[above=7pt,draw=none,font=\small,text=white]{2. autorizar}(s.north west);
\draw[-{Stealth},thick,ouro,bend left=24] (s.south west) to node[below=7pt,draw=none,font=\small,text=white]{3. callback + código}(f.south east);
\end{tikzpicture}
\end{center}
\vfill
{\color{white}\large Exemplo didático com Flask e OAuth 2.0}\par
{\color{white}\large Novembro de 2026}
\end{titlepage}
\pagecolor{white}

\section*{1. Como organizar a integração}
\noindent\colorbox{azulclaro}{\parbox{\dimexpr\linewidth-2\fboxsep}{\textbf{Visão do exemplo.} O tutorial mostra uma implementação de referência em Flask: rotas, serviço OAuth, consulta ao perfil e vínculo local via \texttt{suap\_id}. Adapte os nomes de módulos e configurações ao seu aplicativo.}}
\medskip
O fluxo fica dividido por responsabilidade: as rotas iniciam OAuth e recebem o callback; um serviço monta URLs, troca o código por token e consulta o perfil; outro serviço associa o perfil a um usuário local e chama Flask-Login; o modelo guarda o vínculo; e o template apresenta o botão de entrada. A migration acrescenta a coluna sem apagar os usuários existentes.

\clearpage
\section*{Visão geral do fluxo de autenticação}
O diagrama abaixo mostra o percurso completo da autorização até a sessão Flask-Login. Leia de cima para baixo: o navegador transporta a pessoa entre os sites, enquanto a troca do código por token acontece diretamente entre os servidores.
\begin{center}
\begin{tikzpicture}[x=.9cm,y=.9cm,>=Stealth, every node/.style={font=\large}]
\node[draw=azul,fill=azulclaro,rounded corners=5pt,minimum width=4.3cm,minimum height=1.05cm,align=center] at (0,0) (user) {Pessoa usuária\\\small Navegador};
\node[draw=azul,fill=azulclaro,rounded corners=5pt,minimum width=4.3cm,minimum height=1.05cm,align=center] at (7,0) (app) {Aplicação Flask\\\small Rotas + Flask-Login};
\node[draw=verde,fill=green!8,rounded corners=5pt,minimum width=4.3cm,minimum height=1.05cm,align=center] at (14,0) (suap) {SUAP\\\small OAuth2 + API};
\draw[dashed,gray,thick] (0,-.6)--(0,-9.6);
\draw[dashed,gray,thick] (7,-.6)--(7,-9.6);
\draw[dashed,gray,thick] (14,-.6)--(14,-9.6);
\draw[-{Stealth},thick,azul] (0,-1.15)--(7,-1.15) node[midway,above,fill=white,align=center,text width=5.5cm,font=\normalsize]{1. Clica em “Entrar com SUAP”};
\draw[-{Stealth},thick,azul] (7,-2.25)--(14,-2.25) node[midway,above,fill=white,align=center,text width=5.7cm,font=\normalsize]{2. Redireciona para autorização\\com client\_id, callback e state};
\draw[-{Stealth},thick,verde] (14,-3.35)--(0,-3.35) node[midway,above,fill=white,align=center,text width=6cm,font=\normalsize]{3. SUAP exibe login e consentimento};
\draw[-{Stealth},thick,verde] (0,-4.45)--(14,-4.45) node[midway,above,fill=white,align=center,text width=6cm,font=\normalsize]{4. Pessoa informa a senha no SUAP\\(ela não passa pelo Flask)};
\draw[-{Stealth},thick,verde] (14,-5.55)--(7,-5.55) node[midway,above,fill=white,align=center,text width=5.7cm,font=\normalsize]{5. SUAP redireciona navegador\\ao callback com code + state};
\draw[-{Stealth},thick,ouro] (7,-6.65)--(14,-6.65) node[midway,above,fill=white,align=center,text width=5.7cm,font=\normalsize]{6. Flask troca code por access token\\(client secret fica no servidor)};
\draw[-{Stealth},thick,ouro] (7,-7.75)--(14,-7.75) node[midway,above,fill=white,align=center,text width=5.7cm,font=\normalsize]{7. Flask consulta API com Bearer token};
\draw[-{Stealth},thick,azul] (7,-8.85)--(0,-8.85) node[midway,above,fill=white,align=center,text width=5.7cm,font=\normalsize]{8. Cria/carrega Usuario e abre sessão};
\end{tikzpicture}
\end{center}
\begin{center}\colorbox{azulclaro}{\parbox{.9\linewidth}{\textbf{Leitura rápida:} a senha do SUAP é digitada somente no site do SUAP. O Flask recebe um código temporário; o access token e o client secret são tratados no back-end.}}\end{center}
\begin{center}\small Figura 1 — Fluxo OAuth2 Authorization Code integrado à sessão local do Flask.\end{center}

\section*{2. Preparação: conta, URLs e tipo de cliente}
É necessário ter uma conta ativa no SUAP/IFRN com permissão para registrar aplicações. Acesse \url{https://suap.ifrn.edu.br/api/} e crie uma aplicação OAuth2 (o painel pode pedir autenticação). Use estes valores para uma aplicação Flask executada localmente:
\begin{itemize}[nosep]
\item Nome: por exemplo, \texttt{Minha aplicação Flask (desenvolvimento)}.
\item Tipo de cliente: \textbf{Confidential}, pois o segredo ficará no servidor Flask.
\item Tipo de concessão (grant): \textbf{Authorization Code}.
\item URI de redirecionamento: \texttt{http://127.0.0.1:5000/suap/callback}.
\item Ativo: habilitado; suporte OIDC/algoritmo, se exibido, deve seguir a configuração suportada pelo SUAP e não é necessário para OAuth2 básico.
\end{itemize}
Cadastre a URI exatamente como será usada no Flask: host, porta, caminho e barra final precisam coincidir. Se testar com \texttt{localhost}, cadastre essa URI também; não presuma equivalência entre \texttt{localhost} e \texttt{127.0.0.1}. Para hospedagem, acrescente a URL HTTPS pública do callback. Anote \texttt{Client ID} e \texttt{Client Secret} no momento da criação; trate o segredo como senha e não o coloque no Git.

\clearpage
\section*{Cadastro no painel SUAP: sequência visual}
Use esta sequência junto à tela de cadastro em \url{https://suap.ifrn.edu.br/api/}. Os nomes dos campos podem variar um pouco; os valores do exemplo devem corresponder ao ambiente em que o Flask está rodando.
\begin{center}
\begin{tikzpicture}[>=Stealth,node distance=1.45cm,every node/.style={font=\large}]
\node[draw=azul,fill=azulclaro,rounded corners=6pt,minimum width=14.5cm,minimum height=1.35cm,align=left,text width=13.6cm] (n1) {\textbf{1. Abra o cadastro de aplicações}\hspace{.5em}Entre no SUAP com uma conta autorizada e acesse a área de API / aplicações OAuth2.};
\node[draw=azul,fill=azulclaro,rounded corners=6pt,minimum width=14.5cm,minimum height=1.35cm,align=left,text width=13.6cm,below=of n1] (n2) {\textbf{2. Identifique sua aplicação}\hspace{.5em}Informe um nome reconhecível, por exemplo \texttt{Minha aplicação Flask — desenvolvimento}.};
\node[draw=verde,fill=green!8,rounded corners=6pt,minimum width=14.5cm,minimum height=1.35cm,align=left,text width=13.6cm,below=of n2] (n3) {\textbf{3. Escolha o fluxo para servidor}\hspace{.5em}Client Type: \textbf{Confidential}. Grant Type: \textbf{Authorization Code}. Mantenha a aplicação ativa.};
\node[draw=verde,fill=green!8,rounded corners=6pt,minimum width=14.5cm,minimum height=1.35cm,align=left,text width=13.6cm,below=of n3] (n4) {\textbf{4. Registre o callback exato}\hspace{.5em}\texttt{http://127.0.0.1:5000/suap/callback}. Cadastre também a URL HTTPS de produção quando publicar.};
\node[draw=ouro,fill=yellow!15,rounded corners=6pt,minimum width=14.5cm,minimum height=1.35cm,align=left,text width=13.6cm,below=of n4] (n5) {\textbf{5. Guarde as credenciais}\hspace{.5em}Copie Client ID e Client Secret para variáveis de ambiente; não coloque o segredo no código nem no Git.};
\draw[-{Stealth},very thick,azul] (n1)--(n2);
\draw[-{Stealth},very thick,azul] (n2)--(n3);
\draw[-{Stealth},very thick,azul] (n3)--(n4);
\draw[-{Stealth},very thick,azul] (n4)--(n5);
\end{tikzpicture}
\end{center}
\begin{center}\colorbox{yellow!15}{\parbox{.9\linewidth}{\textbf{Atenção à URI:} host, protocolo, porta, caminho e barra final precisam coincidir literalmente. \texttt{localhost} e \texttt{127.0.0.1} são endereços diferentes para a validação do callback.}}\end{center}
\begin{center}\small Figura 2 — Cadastro de um cliente OAuth2 confidencial para desenvolvimento local.\end{center}

\section*{3. Configurar credenciais e preparar o banco}
Esta etapa parte de um clone limpo. Instale Python 3.12 ou superior e MySQL Server. Se você já recebeu a pasta do projeto, comece na raiz \texttt{app\_flask}. No comando abaixo, substitua o endereço de exemplo pelo endereço Git fornecido pela turma:
\begin{lstlisting}[language=bash]
git clone URL_DO_REPOSITORIO app_flask
cd app_flask
python3 -m venv .venv
source .venv/bin/activate
python -m pip install --upgrade pip
pip install -r requirements.txt
\end{lstlisting}
No Windows, crie o ambiente virtual:
\begin{lstlisting}[language=bash]
py -m venv .venv
\end{lstlisting}
No PowerShell, ative-o e instale as dependências:
\begin{lstlisting}[language=bash]
.\.venv\Scripts\Activate.ps1
py -m pip install --upgrade pip
pip install -r requirements.txt
\end{lstlisting}
No Prompt de Comando (cmd.exe), use:
\begin{lstlisting}[language=bash]
.venv\Scripts\activate.bat
py -m pip install --upgrade pip
pip install -r requirements.txt
\end{lstlisting}
O serviço SUAP usa \texttt{urllib}, biblioteca padrão do Python. Não é necessário instalar \texttt{requests} nem uma dependência OAuth adicional.

O pacote \texttt{mysql-connector-python} instalado pelo pip é o driver de conexão; ele não instala o servidor MySQL. Instale e inicie o MySQL Server no sistema operacional antes de criar o banco. Escolha um nome de banco e use o mesmo valor na configuração \texttt{DB\_NAME} do aplicativo. Neste exemplo, o banco se chama \texttt{suap\_flask\_app}:
\begin{lstlisting}[language=bash]
mysql -u root -p
\end{lstlisting}
No prompt do MySQL, execute:
\begin{lstlisting}[language=SQL]
CREATE DATABASE `suap_flask_app` CHARACTER SET utf8mb4;
EXIT;
\end{lstlisting}
O projeto carrega \texttt{.env} com \texttt{load\_dotenv()} em \texttt{config.py}; o \texttt{.gitignore} já exclui esse arquivo. Crie e edite-o na raiz:
\begin{lstlisting}[language=bash]
nano .env
\end{lstlisting}
No Windows, abra o arquivo com \texttt{notepad .env}. Gere uma chave secreta de sessão:
\begin{lstlisting}[language=bash]
python -c "import secrets; print(secrets.token_hex(32))"
\end{lstlisting}
Copie o resultado e as credenciais da aplicação registrada no SUAP para o arquivo \texttt{.env}. Ajuste usuário e senha de MySQL se necessário:
\begin{lstlisting}
SECRET_KEY=cole_a_chave_gerada
DB_USERNAME=root
DB_PASSWORD=sua_senha_mysql
SUAP_URL=https://suap.ifrn.edu.br
SUAP_CLIENT_ID=cole_aqui_o_client_id
SUAP_CLIENT_SECRET=cole_aqui_o_client_secret
\end{lstlisting}
As configurações correspondentes em \texttt{config.py} são:
\begin{lstlisting}[language=Python]
SECRET_KEY = os.getenv('SECRET_KEY', 'chave-secreta')
DB_PASSWORD = os.getenv("DB_PASSWORD", "labinfo")
DB_USERNAME = os.getenv("DB_USERNAME", "root")
SUAP_URL = os.getenv('SUAP_URL', 'https://suap.ifrn.edu.br')
SUAP_CLIENT_ID = os.getenv('SUAP_CLIENT_ID')
SUAP_CLIENT_SECRET = os.getenv('SUAP_CLIENT_SECRET')
\end{lstlisting}
No arquivo \texttt{app/models/usuario.py}, dentro da classe \texttt{Usuario}, adicione o campo abaixo. Ele é opcional para preservar as contas locais anteriores:
\begin{lstlisting}[language=Python]
suap_id = db.Column(db.String(64), unique=True, nullable=True)
\end{lstlisting}
Com MySQL ativo, o banco criado e as credenciais em \texttt{.env}, crie a migration para registrar a alteração do modelo e aplique-a:
\begin{lstlisting}[language=bash]
flask db migrate -m "adiciona suap_id ao usuario"
flask db upgrade
flask db current
\end{lstlisting}
Abra o arquivo criado em \texttt{migrations/versions/} e confirme que ele adiciona \texttt{suap\_id} como anulável e com restrição única. Neste branch atualizado, a migration já existe: execute somente \texttt{flask db upgrade}, sem gerar outra migration. A pasta \texttt{migrations/} já existe, portanto não execute \texttt{flask db init}. Em um banco vazio, \texttt{upgrade} também cria as tabelas iniciais de usuário e postagem.
\texttt{suap\_id} fica nulo nas contas locais já existentes. Ao entrar pelo SUAP, o serviço associa a conta pelo identificador informado no perfil. A chave gerada e salva como \texttt{SECRET\_KEY} no \texttt{.env} substitui o valor de desenvolvimento padrão de \texttt{config.py}; não exponha o valor padrão em produção.

\section*{4. Como o serviço conversa com o SUAP}
O serviço do projeto está em \path{app/services/SuapAuthenticationService.py}. Ele monta a URL de autorização, troca o código OAuth por token usando \texttt{client\_id} e \texttt{client\_secret}, e consulta o perfil com Bearer token. Abaixo está o código atual:
Se estiver começando a partir de uma versão do projeto sem esse serviço, crie o arquivo no local indicado antes de copiar o código:
\begin{lstlisting}[language=bash]
touch app/services/SuapAuthenticationService.py
\end{lstlisting}
\begin{lstlisting}[language=Python]
import json
from urllib.error import URLError
from urllib.parse import urlencode
from urllib.request import Request, urlopen

from app import app


class SuapAuthenticationService:
    @staticmethod
    def authorization_url(state, redirect_uri):
        params = urlencode({
            'response_type': 'code',
            'client_id': app.config['SUAP_CLIENT_ID'],
            'redirect_uri': redirect_uri,
            'state': state,
        })
        return f"{app.config['SUAP_URL']}/o/authorize/?{params}"

    @staticmethod
    def exchange_code(code, redirect_uri):
        data = urlencode({
            'grant_type': 'authorization_code',
            'code': code,
            'redirect_uri': redirect_uri,
            'client_id': app.config['SUAP_CLIENT_ID'],
            'client_secret': app.config['SUAP_CLIENT_SECRET'],
        }).encode()
        try:
            with urlopen(Request(
                f"{app.config['SUAP_URL']}/o/token/",
                data=data,
                headers={'Content-Type': 'application/x-www-form-urlencoded'},
            )) as response:
                return json.load(response)
        except URLError:
            return None

    @staticmethod
    def meus_dados(access_token):
        request = Request(
            f"{app.config['SUAP_URL']}/api/rh/meus-dados/",
            headers={
                'Authorization': f'Bearer {access_token}',
                'Accept': 'application/json',
            },
        )
        try:
            with urlopen(request) as response:
                return json.load(response)
        except URLError:
            return None
\end{lstlisting}
Os caminhos usados pelo projeto são \path{/o/authorize/}, \path{/o/token/} e \path{/api/rh/meus-dados/}. Para outra instalação SUAP, altere \texttt{SUAP\_URL} e confira os campos do perfil na documentação interativa daquele ambiente.

\section*{5. Rotas Flask: iniciar, validar e concluir o login}
O link do botão em \texttt{app/templates/login.html} aponta para a rota \texttt{login\_suap}:
\begin{lstlisting}[language=HTML]
<a class="btn btn-secondary btn-user btn-block" href="{{ url_for('login_suap') }}">
    Entrar com SUAP
</a>
\end{lstlisting}
Em \texttt{app/routes.py}, a rota de início cria um \texttt{state} aleatório, guarda-o na sessão e encaminha o navegador ao SUAP. A URI de callback é montada a partir do host e da rota atual; ela precisa ser idêntica à cadastrada no painel:
As importações abaixo já existem no arquivo atual e são necessárias para as rotas apresentadas:
\begin{lstlisting}[language=Python]
from app import app
from flask import render_template, redirect, flash, request, url_for, session
from secrets import token_urlsafe
from urllib.parse import urljoin
from app.services.AuthenticationService import AuthenticationService
from app.services.SuapAuthenticationService import SuapAuthenticationService
\end{lstlisting}
\begin{lstlisting}[language=Python]
@app.route('/login-suap')
def login_suap():
    state = token_urlsafe(32)
    redirect_uri = request.host_url.rstrip('/') + url_for('suap_callback')
    session['suap_oauth_state'] = state
    return redirect(SuapAuthenticationService.authorization_url(state, redirect_uri))
\end{lstlisting}
O callback valida o \texttt{state}, troca o código, busca os dados do perfil, vincula o usuário e exibe a página \texttt{suap\_eu.html}. O token fica em memória apenas durante a chamada e não é gravado no cookie de sessão:
\begin{lstlisting}[language=Python]
@app.route('/suap/callback')
def suap_callback():
    state = session.pop('suap_oauth_state', None)
    if request.args.get('error') or not state or state != request.args.get('state', ''):
        flash('Falha na autenticação do SUAP.', 'error')
        return redirect(url_for('login'))

    code = request.args.get('code')
    if not code:
        flash('Código de autorização do SUAP não recebido.', 'error')
        return redirect(url_for('login'))

    redirect_uri = request.host_url.rstrip('/') + url_for('suap_callback')
    tokens = SuapAuthenticationService.exchange_code(code, redirect_uri)
    if not tokens or not tokens.get('access_token'):
        flash('Não foi possível obter o token do SUAP.', 'error')
        return redirect(url_for('login'))

    dados = SuapAuthenticationService.meus_dados(tokens['access_token'])
    if dados is None:
        flash('Não foi possível consultar os dados do SUAP.', 'error')
        return redirect(url_for('login'))

    usuario = AuthenticationService.login_suap(dados)
    if usuario is None:
        flash('Não foi possível vincular esta conta SUAP a um usuário local.', 'error')
        return redirect(url_for('login'))

    foto = dados.get('url_foto_75x100') or dados.get('url_foto') or dados.get('foto')
    foto_url = urljoin(app.config['SUAP_URL'] + '/', foto) if foto else None
    return render_template('suap_eu.html', title='Perfil SUAP', dados=dados, foto_url=foto_url)
\end{lstlisting}
O \texttt{state} vincula o início ao retorno OAuth e evita aceitar callbacks iniciados fora da sessão. Se o SUAP recusar o acesso, retornar sem código ou falhar ao consultar a API, a rota informa o erro e volta à tela de login. Depois do login, \texttt{suap\_eu.html} mostra nome, vínculo, matrícula, e-mail, campus e foto quando esses campos estiverem disponíveis.

\section*{6. Associação do perfil e persistência em \texttt{suap\_id}}
O usuário local continua sendo um modelo \texttt{Usuario} do SQLAlchemy e o Flask-Login mantém a sessão. Em \texttt{app/services/AuthenticationService.py}, o valor de \texttt{suap\_id} vem de matrícula, username ou e-mail do SUAP; o username local vem de nome usual ou nome. O serviço primeiro procura por \texttt{suap\_id}; se ainda não houver vínculo, procura uma conta local compatível e associa-a:
\begin{lstlisting}[language=Python]
from app import db
from app.models.usuario import Usuario
from flask_login import login_user, logout_user
from werkzeug.security import check_password_hash


class AuthenticationService():
    def login(form):
        username = form.username.data.strip()
        user = Usuario.query.filter_by(username=username).first()

        if user and user.password_hash:
            if check_password_hash(user.password_hash, form.password.data):
                login_user(user, remember=form.remember_me.data)
                return True
        return False

    def logout():
        return logout_user()

    def login_suap(dados):
        suap_id = str(
            dados.get('matricula') or dados.get('username') or dados.get('email') or ''
        ).strip()
        if not suap_id:
            return None

        username = str(dados.get('nome_usual') or dados.get('nome') or '').strip()
        email = dados.get('email') or f'{username}@suap.local'
        user = Usuario.query.filter_by(suap_id=suap_id).first()

        if not user:
            user = Usuario.query.filter(
                (Usuario.username == username) | (Usuario.email == email)
            ).first()

        if user and user.suap_id and user.suap_id != suap_id:
            return None

        if not user:
            user = Usuario(username=username, email=email, suap_id=suap_id)
            db.session.add(user)
        else:
            user.suap_id = suap_id

        try:
            db.session.commit()
        except Exception:
            db.session.rollback()
            return None

        login_user(user)
        return user
\end{lstlisting}
O modelo declara o campo assim:
\begin{lstlisting}[language=Python]
suap_id = db.Column(db.String(64), unique=True, nullable=True)
\end{lstlisting}
O campo aceita \texttt{NULL} para não exigir uma identidade SUAP das contas criadas pela aplicação. A restrição \texttt{unique=True} evita associar o mesmo identificador remoto a mais de uma conta. A senha SUAP não é copiada para \texttt{password\_hash}; a autenticação remota é validada pelo protocolo OAuth e a sessão local é aberta por \texttt{login\_user(user)}.

A migration deve ser gerada por cada aluno a partir do modelo atualizado. Execute os comandos a partir da raiz do projeto:
\begin{lstlisting}[language=bash]
flask db migrate -m "adiciona suap_id ao usuario"
flask db upgrade
\end{lstlisting}
O Alembic cria automaticamente um identificador de revisão para o arquivo. Esse identificador pode variar entre cópias; não o escreva manualmente nem copie o de outra pessoa. Antes de aplicar a migration, confira se ela adiciona uma coluna anulável e única. Neste checkout, a migration já está incluída, então basta executar \texttt{flask db upgrade}.

\section*{7. Executar e percorrer o fluxo completo}
Ative o ambiente virtual e instale as dependências declaradas no aplicativo. Confirme que o MySQL está ativo, que o banco escolhido existe e que o arquivo \texttt{.env} contém as credenciais da aplicação. Aplique a migration adicionada e inicie o Flask na mesma porta e host usados no cadastro SUAP:
\begin{lstlisting}[language=bash]
flask db upgrade
flask run --host 127.0.0.1 --port 5000 --debug
\end{lstlisting}
Abra \url{http://127.0.0.1:5000/login}, clique em “Entrar com SUAP”, autentique-se no domínio do SUAP e autorize o acesso. O navegador retorna a \texttt{/suap/callback}; a aplicação valida \texttt{state}, troca o código por token, chama \texttt{/api/rh/meus-dados/}, cria ou encontra o registro com \texttt{suap\_id}, inicia sessão Flask-Login e mostra \texttt{suap\_eu.html}. Confirme o nome e os dados do perfil; em seguida, teste sair por \texttt{/logout} e entrar novamente pelo botão SUAP.

\section*{8. Diagnóstico rápido}
\begin{itemize}
\item \textbf{redirect URI inválida}: o callback usa host da requisição e rota \texttt{suap\_callback}. Compare protocolo, host, porta e caminho com o cadastro no SUAP.
\item \textbf{client inválido}: confira se Client ID/Secret correspondem à aplicação ativa e se o segredo não contém espaços adicionais.
\item \textbf{state inválido}: início e callback precisam compartilhar a mesma sessão/cookie e chave secreta Flask; não inicie o callback diretamente sem começar pelo botão.
\item \textbf{token aceito, perfil negado}: confira o endpoint \texttt{/api/rh/meus-dados/}, as permissões da conta e a documentação interativa do SUAP. O serviço envia o token como \texttt{Authorization: Bearer ...}.
\item \textbf{funciona em localhost, falha publicado}: cadastre a URL HTTPS pública do callback e configure proxy/HTTPS corretamente; a rota deriva o callback de \texttt{request.host\_url}.
\item \textbf{usuário autentica, mas áreas protegidas redirecionam}: confira vínculo ao modelo \texttt{Usuario}, chamada a \texttt{login\_user} e função \texttt{load\_user}.
\end{itemize}

\section*{9. Referências e observação de segurança}
\begin{itemize}
\item Documentação interativa oficial: \url{https://suap.ifrn.edu.br/api/docs/}.
\item Cadastro e acesso à API: \url{https://suap.ifrn.edu.br/api/}.
\item Cliente oficial IFRN para acesso à API: \url{https://ifrn.github.io/suapi/}.
\item Referência IFRN de aplicação OAuth2 Confidential + Authorization Code: \url{https://github.com/ifrn-oficial/cliente_suap_django}.
\end{itemize}
Os formulários e rótulos do painel podem mudar. A documentação interativa do SUAP é a fonte final para endpoints, parâmetros e permissões vigentes. O código consulta o endpoint de perfil configurado no serviço e descarta o access token ao final da requisição; nunca compartilhe Client Secret, access token ou cookies de sessão.
\end{document}
